\section{Introduction}
\label{sec:introduction}

A redistricting analysis that evaluates the enacted map against seven
fairness metrics is simultaneously testing seven hypotheses:
\begin{enumerate}
\item $H_{\PP}$: The plan's Polsby-Popper compactness is within the
  normal range for valid plans.
\item $H_{\rm EG}$: The plan's efficiency gap is within the normal range.
\item $H_{\rm MM}$: The plan's mean-median difference is within the
  normal range.
\item $H_{\rm PB}$: The plan's partisan bias (at the 50\% vote-share
  threshold) is within the normal range.
\item $H_{\rm SV}$: The plan's seats-votes curve slope is within the
  normal range.
\item $H_D$: The plan's Democratic seat count is within the normal range.
\item $H_{\rm DG}$: The plan's declination metric is within the normal range.
\end{enumerate}

If any one of these hypotheses is rejected, the analyst might claim
evidence of a gerrymander.
But testing seven hypotheses simultaneously at $\alpha = 0.05$ creates
a family-wise error rate of $1 - (0.95)^7 \approx 30\%$: a 30\% chance
of falsely concluding that the plan is extreme on at least one metric
even if it was drawn randomly.

This paper addresses the multiple testing problem systematically by
applying the Benjamini-Hochberg (BH) FDR procedure~\citep{bh1995} to
the full L-track test battery.

\subsection{Scale of the Testing Problem}

The L-track's complete test battery for the 2020 redistricting cycle
involves:
\begin{itemize}
\item 7 partisan fairness metrics
\item 50 states (including 7 with $k > 1$ and redistricting obligations)
\item 3 census cycles (2000, 2010, 2020)
\item Total tests: $7 \times 50 \times 3 = 1{,}050$
\end{itemize}

At $\alpha = 0.05$ without correction:
\[
  E[\text{false positives}] = 1{,}050 \times 0.05 = 52.5.
\]
Over 50 tests are expected to show ``significant'' gerrymandering evidence
in a world where no state has been gerrymandered.

\subsection{The Holm Implementation}

\bisect{}'s \texttt{bisect-analysis/src/bloc\_voting.rs} already implements
Holm-Bonferroni correction for bloc voting analysis.
The Rust code computes a \texttt{p\_value\_holm} field alongside the
raw p-value and the Bonferroni-corrected p-value.

The Holm-Bonferroni procedure controls the FWER: the probability that any
single test is a false positive.
This is appropriate for the bloc voting context (where we want to ensure
that no VRA district claim is a false positive) but overly conservative
for the L-track portfolio analysis.

The BH FDR procedure is more appropriate for the portfolio analysis:
it allows a controlled proportion of false positives (FDR $\leq q$)
rather than requiring that the probability of any false positive be
$\leq \alpha$ (FWER control).
When many hypotheses are expected to be true (as in a multi-state
analysis where many states do have partisan maps), FDR control provides
substantially more power than FWER control.

\subsection{Road Map}

Section~\ref{sec:testing} describes the testing problem in detail.
Section~\ref{sec:bh} presents the BH procedure.
Section~\ref{sec:impl} describes the existing Holm implementation and
how S.4 extends it.
Section~\ref{sec:application} applies BH to the L-track battery.
Section~\ref{sec:results} presents the corrected findings.
Section~\ref{sec:conclusion} concludes.
